% ===========================================================================
% Section 12. Discussion, next steps and conclusion
% Numbers: paper/results/*.json as reported in Sections 5-11: the indices, picks and partnership
% index from the final run of 12 October 2026; the Allied Industrial Base from aib.json and
% aib_robustness.json (AIB run of 16 October 2026 after the credibility pass). Conceptual frame
% and reference list: the NEIS working paper "Allied Industrial Base: Conceptual overview,
% literature review and analytical design" (October 2026), cited as neis2026allied; paraphrased,
% not quoted.
% ===========================================================================
\section{Discussion, next steps and conclusion}\label{sec:discussion}

This section reads the results through the analytical design for an Allied Industrial Base of
the New Energy Industrial Strategy working paper \citep{neis2026allied}, states what the paper does
not establish, and sets out the extensions that design calls for.

\subsection{What the results imply for an allied industrial base}\label{subsec:discussion_implications}

\paragraph{Managed interdependence rather than autarky.} The working paper defines the objective
as managed interdependence at allied scale: strategically consequential capabilities should exist
at sufficient scale somewhere in a trusted network, not in every member, because no single ally
can reproduce China's scale across the electro-industrial stack
\citep{neis2026allied,campbell2025underestimating}. China leads 22 of the 33 Economic Opportunity pairs and 11
of the 30 Energy Security pairs (Sections~\ref{sec:es} and~\ref{sec:eo}); a coalition member leads
16 security pairs but only 8 opportunity pairs. Within each member, opportunity and security need
move in opposite directions: the median within-country Spearman correlation between $\EO$ and
$1-\ES$ is $-0.49$ and negative for every member (Section~\ref{sec:strategic}). A member's strongest
opportunities lie where it is already secure, and its needs lie where others are strong. The
Allied Industrial Base analysis shows what this means under stress (Section~\ref{sec:network}). With
China's supply removed, the 57 members lose 22.9 percent of their traded electro-industrial demand.
Existing relationships restore 7.3 points of it and optimized partnerships 0.3 more: the coordination
dividend is below one point in 81 of the 91 nodes. The trade partners the allies need are already
their partners; what is missing is supply.

The allied response is not new supply, either. Under the headline reading of the shock, all of the
recovery, USD 232 billion in the electro-industrial scope, is drawn from buyers outside the coalition,
and with every importer's demand protected nothing can be recovered by trade at all
(Section~\ref{subsec:robust_network}). China is the largest world exporter in 9 of the 18 basket
stages, including eight of the nine Midstream ones, with shares of 19 to 44 percent. The evidence that
broad relocalisation is costly and that diversification usually serves resilience better than
reshoring \citep{oecd2025supply,thakurweigold2024promoting,aiyar2023geoeconomic} applies to an allied
wall as much as to a national one. The stress test adds a caution about diversification itself:
against a China shock a supplier-share limit costs little, but in 63 of the 91 nodes no tested limit
restores 80 percent of demand after a second loss, because the capacity to diversify into does not
exist.

\paragraph{The node, not the sector, as the unit of treatment.} The working paper argues that
policy should be set by node, because one technology contains very different risks
\citep{neis2026allied}. The AIB analysis works at the sub-sector node, and the stages of the indices
show the same heterogeneity. Technology and stage explain 44 percent of the variance of $\ES$ and 58
percent of $\EO$, against 12 and 18 percent for countries alone, and concentration sits in different
places along each chain. Midstream security is mostly a trade question (the Trade category supplies 73
percent of the Midstream score), and China is the largest exporter in almost every Midstream basket.
Upstream supply is led by resource holders instead: Australia in the Batteries, Electric Vehicles and
Coal baskets, Chile in Electric Grid (copper) and Kazakhstan in Nuclear. Chinese-parent ownership of
foreign-invested plants is high in solar (median share 0.75 across partners) and nil in wind
(Section~\ref{sec:bilateral}). Within the electro-industrial chains, 37 of the 70 nodes need new
capacity and 32 need no intervention, so a rule set by technology would treat most nodes wrongly.

Node-level treatment departs from naive rankings in two places. First, the need reading of the
Strategic Index rates a stage highest where a member lacks the resource. Without the endowment
screen 47 of the 57 allied picks would be Upstream, 17 of them in Nuclear Upstream, many in members
with no uranium. With the screen 43 picks are Midstream, and 41 of the 57 picks change (Figure~\ref{fig:strategic_screen}). The
eight Nuclear Upstream picks that remain all hold identified uranium resources, but small ones
(0.015 to 0.79 percent of the world's) and none with recorded mine production in 2019--2024, so
they are still need-driven: the screen removes the no-resource artefact but not its small-resource
counterpart. The screen is a rule of pick eligibility and leaves every index value unchanged. Nor
are the picks an allied assignment: 27 members would prioritise Batteries Midstream, but the minimum
coalitions of the four compacts with a Batteries Midstream node contain 12 of them, and only 13 of the
57 members sit in a coalition for their pick stage. Second, the stress test does not follow the export
ranking either: the economies that would host the most new capacity, Hungary, Czechia, the Philippines
and Viet Nam, are minor suppliers of added trade, and hosts qualify on revealed capability at scale,
not on their rank (Section~\ref{subsec:aib_investment}).

\paragraph{Partners in the emergent network.} The working paper asks
an allied design to ``optimize combinations, not rankings'', valuing a partner for the capability
it adds rather than for its stand-alone score \citep{neis2026allied}. As bilateral partners the 26
partner-tier members rank low: their median rank in the headline Partnership Strength Index is 64.5
of 192, against 16 for the core members, in part because UNGA alignment makes up on average 24
percent of the Friendshore composite. In the stress test they matter more. Partners would host 33.6
percent of the supply from the leading hosts' new capacity and economies outside the coalition 7.4
percent; Viet Nam sits in 16 of the 22 compact coalitions, more than any member; and of the six
critical supply partners, two are partners (Singapore and Viet Nam) and three are outside the
coalition altogether (Madagascar, Mozambique and Taiwan). Economies outside the coalition also provide
14.4 percent of the added trade. The coalition is the set of economies the analysis protects; the
industrial base that supplies it is wider.

\paragraph{The three hypotheses.} The working paper proposes three hypotheses for testing first
\citep{neis2026allied}. The paper tests none of them; its results bear on each to a different
degree. On the first, that allied specialisation reaches a resilience target at lower public cost
than parallel national programmes, the compacts show how few countries a target needs (three to nine
reach 95 percent of the anchors' traded demand in every node), but the analysis contains no costs or
outlays. On the second, that coordinating demand matters more than coordinating subsidies, the results
point to capacity: the trade compacts add at most 0.04 points to the system's coverage, and every
compact is an investment and offtake club, which makes demand aggregation, offtake and finance rather
than trade rules the instruments that matter; whether demand or subsidy coordination delivers the
capacity is beyond the data. On the third, that rules targeted at high-severity nodes deliver most of
the benefit of broad decoupling at lower cost, the results document a precondition, heterogeneity
across nodes and stages (the Strategic Index is 73 percent sector and 4 percent country variance), but
not the benefit or cost of either rule, because the stress test removes China from every node, itself a
sector-wide shock.

\paragraph{Normative choices and how readers can vary them.} The weights of the Strategic Index,
its need reading, the alignment weight in Friendshore and the coalition thresholds are choices, and so
are, in the AIB analysis, the economies protected, the reading of the shock and the classification
thresholds. As the working paper argues, such weights are political, and analysis should display the
frontier of choices rather than one number \citep{neis2026allied}. Some choices matter a great deal
(Section~\ref{sec:robustness}): reading the security term as strength rather than need keeps only 5 of
the 57 picks; random draws of the five Strategic Index weights keep a member's pick in 73 percent of
draws on average; dropping alignment from Friendshore moves China from 44th to 3rd in the Partnership
Strength Index; reading the China shock as a decoupling rather than an export restriction lowers the
electro-industrial nodes that need new capacity from 37 to 24. The interactive explorer that accompanies
the paper exposes the index choices through four Perspective dials in its country profiles:
\emph{hawkishness} weights security (the need term, Friendshore, alignment and Chinese-parent
exposure), \emph{bullishness} opportunity (the $\EO$ term and partnership opportunity),
\emph{populism} policy signals and home manufacturing, and \emph{greenness} climate performance,
dropping coal, oil and gas at its highest settings. The paper's headline is one setting of these dials.
The sector profiles show the theme data (trade flows, reserves, production, investment including the
CM-FDI projects, policy records) by technology and stage without index weights, and the Allied
Industrial Base section shows the stress test, the investor--host flows, the country roles and the
compacts, with observed facts and settings that are preferences labelled apart.

\subsection{Limitations}\label{subsec:discussion_limitations}

\emph{No outcome validation.} The indices are not validated against disruptions, investment or
export growth. The one predictive test, in the study of vertical complementarities
(Section~\ref{subsec:robust_vertical}), is a warning: sourcing-weighted input security did not
predict the import shortfalls of 2021--22, whereas sourcing concentration did. The AIB analysis is
validated only by tracing four nodes through every step (Online Appendix~\ref{oa:network}), not
against a realised shock.

\emph{Data vintages and source conventions.} Inputs come from different years (92 percent of
opportunity cells mix years; fossil reserves date from end-2020, mineral reserves from end-2025),
and the 2020 and 2025 vintages agree less for $\EO$ (median within-pair Spearman 0.87) than for
$\ES$ (0.99). Holdings the USGS does not quantify (not available, withheld or small) count as zero
in the technology roll-ups. The IEA's 2035 announced manufacturing-capacity shares are regional for
Europe, so every EU member carries the EU's share (for example 0.15 in Wind Midstream). The
BloombergNEF outlook behind Downstream Consumption does not report Canada, whose category is imputed
in all eight cells it enters, leaving six of Canada's nine scored Downstream security cells
unrankable.

\emph{Mirror trade.} Economies that do not report to UN Comtrade enter the trade matrix only
through their partners' declarations: 65 of 208 importers in 2024, among them two coalition members
(Gabon and Viet Nam). Electricity is excluded from the grid basket, and Hong Kong's re-exports from
the export measures.

\emph{Traded demand and synthetic capacity.} The AIB analysis covers traded demand only, so
production used at home is outside it; it rests on one reference year, applies each shock node by node
without propagation along the chain, reads trade values as capability except in the 11 nodes with
credible physical data, and adds capacity as a standardised increment rather than as projects with
costs and lead times.

\emph{Policy basis.} The Policy index of 150 economies, 11 of them coalition members, rests on IEA
PAMS alone. For the 65 economies with both components, ranking their cells on IEA PAMS alone rather
than on the composite gives a median within-country Spearman correlation of only 0.48.

\emph{Thin and construction-driven picks.} The pick is within 0.005 of its runner-up for 10 of the
57 members and within 0.01 for 21. The Nuclear Midstream picks of Costa Rica, the Dominican Republic,
Guatemala and Namibia are not an Energy Security floor effect. Each may pick only four sectors, the
Midstream stages of nuclear, oil, gas and coal (Namibia also Nuclear Upstream); Midstream nuclear $\ES$ is compressed around its
median; and the technology weight (0.999 for nuclear against 0.42 for gas, 0.02 for coal and 0 for
oil) decides the pick. The picks reflect a narrow choice set, not an opportunity in nuclear fuel.

\emph{Coalition boundary.} Taiwan (no UN vote, no listed agreement, at most 1.5 percent of any
world stage) and Switzerland (whose only supplier route rested on electricity exports) are outside
the coalition. The AIB analysis does not protect their demand but counts their supply: Taiwan is a
critical supply partner of the battery compact.

\emph{Layers not yet measured.} Of the working paper's five analytical layers, the paper measures
node severity as country-relative exposure and, in the AIB analysis, as the share of traded demand a
shock removes; capability through $\EO$ and revealed export capability; partner complementarity
through the bilateral indices and the compacts; and robustness through the stress tests and an N-1
test. It does not measure the fourth layer, bankability, nor recovery time, common-mode risk beyond
the ownership of foreign-invested plants, or the dynamic option value of latent capacity.

\subsection{Next steps and recommendations}\label{subsec:discussion_next}

The working paper orders the decision in three steps, severity, coordination and specialisation
\citep{neis2026allied}. This paper offers measured versions of the first two for traded goods and a
first version of the third: the compacts name where capacity should be built and with whom.
Table~\ref{tab:discussion_agenda} lists ten extensions drawn from the working paper's
implementation questions, analytical upgrades and research agenda, with the data and model change
each requires. Five deserve emphasis.

First, a \emph{complementarity matrix}. The AIB analysis already prices a country's marginal
contribution to a compact: the leave-one-country-out test and the exact coalition identify the
members no single other economy can replace. Extending the test to every economy and node, with
Shapley values over the coalition program, would replace the stand-alone ranking of partners with the
gap-filling value the working paper asks for.

Second, \emph{common-mode risk and resilience floors}. Two suppliers in two countries do not
diversify supply if they share a parent, a precursor or an equipment maker
\citep{neis2026allied,farrell2019weaponized}. The Chinese-parent component of Friendshore covers
the first channel in four Midstream stages; the trade program could also cap, by constraint, the
supply resting on any one parent headquarters or precursor source. The N-1 test of the stress test is
a first resilience floor. Floors tied to recovery time need capacities and inventories, so that
removing the largest producer for 6, 12 or 24 months yields unmet demand over time. The two other
upgrades need new inputs. \emph{Dynamic option value}, the time and cost of ramping or requalifying
capacity, calls for a multi-period version of the investment step with build lead times, which the
CM-FDI announcement and completion years could calibrate. \emph{Bankability} needs the cost of capital,
permitting times, offtake and delivery capacity by country and stage \citep{spokas2025systemic},
entered as a host-eligibility condition or a cost of the increment.

Third, \emph{chain pilots} for the three chains the working paper proposes: rare earths and magnets,
batteries and anodes, and solar wafers. The AIB analysis already traces permanent magnets, rare-earth
compounds, anode materials and wafers through every step and finds each a capacity gap under the
headline, but each chain must be split into finer sub-stages (separation, alloy, magnet; graphite,
spherical graphite, anode; polysilicon, ingot and wafer, cell) with plant-level capacity and
qualification times. The output would be a portfolio per chain: capacity by host, minimum redundancy,
priority projects and financing gaps.

Fourth, \emph{differentiated Gates and Fences and sectoral compacts}. Combined with severity data
on consequence and substitutability, the node evidence assembled here (exposure, residual gap,
concentration, Chinese-parent exposure, alignment) can sort nodes into high, medium and low fences.
The node table of the compacts hands each node to that evidence with its treatment, and the
investor--host pairs show where development and export finance would back new capacity (Czechia, Viet
Nam and the Philippines lead the most nodes).

Fifth, a \emph{system scorecard with off-ramps}. A recurring rerun already produces several of the
working paper's system metrics by node: exposure, coverage, the residual gap, the largest supplier's
share and the share of supply from inside the coalition. With inventory coverage, recovery time and
the price premium added, it could trigger reviews automatically, for instance when a node's residual
gap rises above the \textsc{build} threshold, and retire support when it falls.

The hypotheses can then be tested. Restricting each member to its own supply gives the
self-sufficiency counterfactual for the first, once costs enter. CM-FDI project histories can be set
against demand instruments and supply subsidies for the second. Shocks that remove China only in
high-severity nodes, rather than in every node, would test the third, and the supplier-share frontier
already shows when another supplier stops adding resilience. The study of vertical complementarities
adds its own agenda (last row of the table).

\begin{table}[htbp]
  \centering
  \begin{threeparttable}
  \caption{Next steps: extensions of the indices and the Allied Industrial Base analysis}
  \label{tab:discussion_agenda}
  \footnotesize
  \begin{tabularx}{\textwidth}{@{}>{\raggedright\arraybackslash}p{0.2\textwidth}
      >{\raggedright\arraybackslash}X >{\raggedright\arraybackslash}X@{}}
    \toprule
    Extension & Data needed & Model change \\
    \midrule
    Complementarity matrix & None beyond the present inputs & Leave-one-country-out over every node and economy; Shapley values over the coalition program \\
    Common-mode risk & Plant ownership beyond foreign investment; precursor and equipment sourcing by plant & Constraint on the supply sharing a parent headquarters or precursor source \\
    Dynamic option value & Project lead times (CM-FDI announcement and completion years), pipeline capacity, technology readiness of substitutes & Multi-period investment step with build lead times \\
    Bankability and feasibility & Cost of capital, permitting times, offtake, insurance, delivery capacity \citep{spokas2025systemic} & Bankability as a host-eligibility condition or a cost of the increment \\
    Resilience floors and recovery time & Demand volumes for 2030 and 2035, plant capacities, inventories & Capacities and inventories in the stress test; outages of 6, 12 or 24 months \\
    Chain pilots & Sub-stage HS6 codes, plant-level capacity, qualification times for magnets, anodes and wafers & Sub-stage nodes for three chains; portfolio by chain \\
    Gates and Fences by node & Consequence and substitutability by node; Chinese-parent exposure beyond four stages & Node-specific access rules for non-members \\
    Sectoral compacts & Offtake and procurement commitments; development and export finance & Offtake volumes and finance by compact; costs of the investment step \\
    System scorecard and off-ramps & Inventory coverage, price premium, capital mobilised & Annual rerun with review triggers on exposure and residual gaps \\
    Vertical complementarities & Plant-level co-location; input--output coefficients; 2023 graphite and gallium controls & Coupling of supply across the nodes of a chain; tested before adoption \\
    \bottomrule
  \end{tabularx}
  \begin{tablenotes}[flushleft]\footnotesize
    \item \emph{Notes:} Extensions drawn from the implementation questions, analytical upgrades and research
    agenda of the NEIS working paper on the Allied Industrial Base \citep{neis2026allied}, matched to the
    data and model of this paper. Model changes extend the headline configuration of
    Section~\ref{sec:network} (China removed, hold, production-constrained, reallocation).
  \end{tablenotes}
  \end{threeparttable}
\end{table}

\subsection{Conclusion}\label{subsec:conclusion}

This paper has built indices of energy security and economic opportunity for 219 economies, eleven
technologies and three supply-chain stages; derived a coalition of 57 members from UN voting and
supply shares; scored every ordered pair of economies as friendshoring opportunities and export
markets; and stress-tested the allied supply network at 92 sub-sector nodes, from the shock through
existing and optimized trade to capacity and Sectoral Compacts. The results support the working
paper's premise in a limited sense: no member is both secure and competitive across the chains, and
the allies' exposure to a China shock is a gap in capacity that trade among them cannot close.
Whether the compacts are cheaper, more bankable or faster to recover than national programmes needs
costs, volumes and time. The indices and the AIB analysis are best read as the analytical backbone of
a recurring review, with the explorer showing how its answers change when the settings change.
